\section{Introduction}

Let $v_1,\ldots,v_N$ be vectors in the Euclidean unit ball of $\R^d$.
There are two natural ways to keep their partial sums small: choose signs
in a prescribed order, or, when the family sums to zero, choose an order
without changing any signs. The difficulty in the first problem is to
control every prefix by one choice of signs, with a bound independent of
$N$. We prove such a bound on the square-root scale in the dimension.

\begin{theorem}[Prescribed-order signed prefixes]\label{thm:signing}
There is an absolute constant $C$ such that, for all integers $d,N\ge1$ and
every sequence $v_1,\ldots,v_N\in\R^d$ with $\lVert v_i\rVert_2\le1$,
there are signs $\varepsilon_1,\ldots,\varepsilon_N\in\{-1,1\}$ for which
\[
 \max_{0\le k\le N}
 \left\lVert\sum_{i=1}^k\varepsilon_i v_i\right\rVert_2\le C\sqrt d.
\]
The same signs control all prefixes, and $C$ is independent of $d$, $N$
and the sequence.
\end{theorem}

The zero-sum ordering problem is quantified by the Euclidean Steinitz
constant. For an indexed family with $\sum_{i=1}^N v_i=0$, define
\[
 \beta(v_1,\ldots,v_N)
 =\min_{\pi\in\mathfrak S_N}\;
   \max_{0\le k\le N}\left\lVert\sum_{i=1}^k v_{\pi(i)}\right\rVert_2,
\]
where $\mathfrak S_N$ is the set of permutations of $\{1,\ldots,N\}$
and an empty sum is zero. Let $S_2(d)$ be the supremum of $\beta$ over
all finite indexed zero-sum families of Euclidean unit-ball vectors in
$\R^d$. The Euclidean Steinitz--Bergstr\"om conjecture asks whether
$S_2(d)=O(\sqrt d)$. Theorem~\ref{thm:signing} gives its affirmative resolution.

\begin{theorem}[Euclidean Steinitz--Bergstr\"om bound]\label{thm:steinitz}
There is an absolute constant $C$ such that, for all integers $d,N\ge1$ and
every indexed family $v_1,\ldots,v_N\in\R^d$ satisfying
$\lVert v_i\rVert_2\le1$ and $\sum_{i=1}^N v_i=0$, some permutation
$\pi\in\mathfrak S_N$ satisfies
\[
 \max_{0\le k\le N}\left\lVert\sum_{i=1}^k v_{\pi(i)}\right\rVert_2
 \le C\sqrt d.
\]
The constant in Theorem~\ref{thm:signing} also suffices here.
\end{theorem}

Both statements allow repetitions and zero vectors. The empty-sequence
and dimension-zero extensions have every partial sum equal to zero. The construction is
existential: its choices may depend on the complete input sequence. Its
successive choice of signs does not supply an online rule or an efficient
algorithm.

The order of growth in Theorem~\ref{thm:steinitz} is necessary. In the
$d$-dimensional hyperplane orthogonal to the all-ones vector
$\mathbf1\in\R^{d+1}$, let $e_i$ denote the standard coordinate vectors and take
\[
 u_i=\sqrt{\frac{d+1}{d}}\left(e_i-\frac{\mathbf1}{d+1}\right),
 \qquad 1\le i\le d+1.
\]
These unit vectors sum to zero and satisfy $\langle u_i,u_j\rangle=-1/d$
for $i\ne j$. Hence every sum of $k$ distinct members has squared norm
$k(d+1-k)/d$. With $k=\lfloor(d+1)/2\rfloor$, this is at least $d/4$,
regardless of the ordering. Thus
\[
 \tfrac12\sqrt d\le S_2(d)\le C\sqrt d\qquad(d\ge1).
\]

\subsection{From signs to an ordering}

We give the short transference argument at once. Chobanyan's
transference principle~\cite{Chobanyan1994} relates rearrangements and
signed sums. We use its finite positive-forward, negative-reverse form
\cite[Theorem~2.1 and Remark~1]{ChobanyanEtAl2023}.

\begin{proof}[Proof of Theorem~\ref{thm:steinitz} from Theorem~\ref{thm:signing}]
Fix an indexed zero-sum family and choose an order attaining its minimum
$\beta$; this exists because there are only finitely many permutations.
Relabel the vectors in that order. Put
\[
 A_j=\sum_{i=1}^jv_i,
 \qquad B_j=\sum_{i=1}^j\varepsilon_i v_i,
 \qquad 0\le j\le N,
\]
where the signs come from Theorem~\ref{thm:signing}. Hence
$\lVert A_j\rVert_2\le\beta$ and $\lVert B_j\rVert_2\le C\sqrt d$.
The positive-sign and negative-sign selections from an original prefix
sum respectively to
\[
 P_j=\frac{A_j+B_j}{2},\qquad M_j=\frac{A_j-B_j}{2}.
\]
Place all positive-sign indices in their original order, followed by all
negative-sign indices in reverse original order. A partial sum in the
first block is some $P_j$. In the second block, the unused indices are
the negative-sign indices of an original prefix. Since the total sum
is zero, the current partial sum is therefore some $-M_j$. This also
covers either block being empty. Every new partial sum has norm at most
$(\beta+C\sqrt d)/2$. By minimality,
\[
 \beta\le\frac{\beta+C\sqrt d}{2},
 \qquad\text{and hence}\qquad \beta\le C\sqrt d.
\]
This permutes indices and so preserves all multiplicities.
\end{proof}

The signed-prefix bound also yields finite-dimensional $\ell_p$ estimates
and a colorful row-permutation theorem. We develop these distinct
consequences in Section~\ref{sec:consequences}.

\subsection{Historical context}

The finite rearrangement theorem goes back to Steinitz's work on
conditionally convergent vector series; Ambrus and Heck give a historical
account in \cite[Section~2]{AmbrusHeck2026}.
The dimension bound for zero-sum rearrangement holds for arbitrary norms:
Grinberg and Sevast'yanov~\cite[Theorem~1]{GrinbergSevastyanov1980} prove
that the partial sums can be kept within $d$ times the unit ball.
The smaller square-root scale in Euclidean space has a long history;
Behrend discusses this expected growth in
\cite[p.~108]{Behrend1954}. The modern formulation and its attribution
to Bergstr\"om are recorded by Ambrus and Heck
\cite[Conjecture~5]{AmbrusHeck2026}. Their Theorem~7 reduces the
asymptotic question to a relaxed rearrangement problem for vectors with
norms in $[1-\varepsilon,1]$, with $0<\varepsilon\le1$ fixed independently
of the dimension. In that problem each prefix is compared with its
proportional share of the total sum, rather than with zero.
For prescribed-order Euclidean signing, the square-root-dimension
question is stated explicitly by Bansal, Jiang, Meka, Singla and
Sinha~\cite[Section~6, Conjecture~6.3]{BansalEtAl2021Full}, following
Banaszczyk~\cite{Banaszczyk2012}. Their discussion records Banaszczyk's
bound $O(\sqrt d+\sqrt{\log N})$; the length term becomes significant when
$N$ is large compared with the dimension. In a 2026 preprint, Dutta, Jha
and Jiang state an efficient prescribed-order signing bound
\[
 O\bigl(\sqrt d+d^{1/4}\log^{7/4}N\bigr),
\]
and the corresponding zero-sum ordering bound
\cite[Theorem~1.1 and Corollary~1.2]{DuttaJhaJiang2026}.
Theorem~\ref{thm:signing} gives an absolute multiple of $\sqrt d$
for every finite length.

Terminal vector balancing has a different quantifier. For columns in the
Euclidean unit ball, Li~\cite[Theorem~1.1]{Li2026} bounds the final
signed sum in infinity norm by an absolute constant and gives a
deterministic polynomial-bit algorithm on rational input. That terminal
statement implies a Euclidean bound of order $\sqrt d$ for one sum.
Applying it separately to each prefix permits the signs to change with
the prefix; Theorem~\ref{thm:signing} instead supplies one common signing.
The present all-prefix estimate is Euclidean and does not assert a
constant infinity-norm all-prefix bound.

\subsection{Proof strategy}

Gaussian convex bodies are a central setting for vector balancing, notably
in the theorem of Banaszczyk~\cite{Banaszczyk1998}. Our proof takes place
in a space of coefficients: one coordinate is
reserved for each incoming vector, together with $d$ initial state
coordinates. A sequence of linear maps sends these coefficients to
states in $\R^d$. At each step an invertible contraction acts on the old
state and the next coefficient injects a multiple of the incoming vector.
The body of coefficients for which every state is bounded will provide
the simultaneous control of all prefixes.

Two analytic ingredients make this construction possible. First, if a
bounded symmetric convex body has small Dirichlet energy for every
positive diagonal weight, it supports a single square root of a probability density
with small derivatives in every coordinate. A lifted Steiner
symmetrization preserves these bounds and permits a coordinate step
of the form $q+\varepsilon$, where $|q|\le1/4$ and
$\varepsilon\in\{-1,1\}$. This adapts the directional-density and
symmetrization method of Guo, Fang and Lu
\cite[Sections~3--4]{GuoFangLu2026}.\footnote{Guo, Fang and Lu
credit the Odin Automatic AI Research Agent for their proof.}
Bandeira~\cite{Bandeira2026} gives a related formulation for terminal
balancing in terms of Dirichlet eigenvalues. Akbas and Sra
\cite[Lemmas~2.2--2.3]{AkbasSra2026} develop the common quadratic-energy
witness and lifted $L^2$-symmetrization argument. Here the hypotheses use
diagonal weights, and the coordinate shifts may depend on earlier choices.
We prove the required energy and shifted-step statements in full.

Second, the filtered body has a diagonal Dirichlet-energy bound
independent of the number of filters. To prove this, we follow a stationary Ornstein--Uhlenbeck process
in coefficient space and localize its filtered states to dyadic spectral
scales. A matrix estimate controls changes at one scale by an additive
trace budget, including gains inserted between noncommuting contractions.
At a fixed scale, this estimate gives a uniform positive probability
for all constraints in a group of indices with bounded accumulated energy,
over a matched time interval. Gaussian correlation combines these groups,
scales and time intervals; the exponential survival rate yields the energy
bound.

The two ingredients meet through a scalar predictor. The allowed shift
in the incoming coefficient exactly compensates for the contraction of
the old state. Consequently the final states are the signed partial
sums themselves. We give this complete construction in
Section~\ref{sec:construction}, after stating the two analytic
propositions. Section~\ref{sec:shifted-signs} proves the shifted-sign proposition.
Section~\ref{sec:matrix} establishes the matrix variation estimate, and
Section~\ref{sec:survival} uses it to prove the filtered-body energy bound.
Section~\ref{sec:consequences} develops the norm-comparison and colorful
row-permutation consequences.

Throughout, all unspecified constants are absolute. Vectors and matrices
are real; $\lVert\cdot\rVert_\op$ and $\lVert\cdot\rVert_\HS$ denote
operator and Hilbert--Schmidt norms. For symmetric matrices, $A\le B$
means that $B-A$ is positive semidefinite. Functions of positive definite
matrices are defined by the spectral calculus.
